% VI33 detailed exercise solutions -- VI/24+ proof-depth standard.

% VI33-DETAILED-EXERCISE-01
\paragraph{Solution VI.33.E1.}
For every open \(U\subseteq X\),
\[
\mathcal O_X^{\oplus r}(U)=\mathcal O_X(U)^{\oplus r},
\]
with scalar multiplication and addition defined componentwise.  If \(V\subseteq U\), the restriction map is the
direct sum of \(r\) copies of \(\mathcal O_X(U)\to\mathcal O_X(V)\), so it is linear with respect to restriction of
scalars. Hence this is an \(\mathcal O_X\)-module.

Stalks commute with finite direct sums because filtered colimits do:
\[
(\mathcal O_X^{\oplus r})_x
=
\varinjlim_{x\in U}\mathcal O_X(U)^{\oplus r}
\cong
\left(\varinjlim_{x\in U}\mathcal O_X(U)\right)^{\oplus r}.
\]
Therefore
\[
\boxed{(\mathcal O_X^{\oplus r})_x\cong\mathcal O_{X,x}^{\oplus r}}.
\]

% VI33-DETAILED-EXERCISE-02
\paragraph{Solution VI.33.E2.}
For \(X=\Spec A\), the structure sheaf is the affine module sheaf associated with the \(A\)-module \(A\):
\[
\mathcal O_X\cong\widetilde A.
\]
By the defining distinguished-open formula,
\[
\widetilde A(D(f))=A_f.
\]
Thus
\[
\boxed{\widetilde A(D(f))\cong A_f}.
\]
Under this identification, restriction \(D(g)\subseteq D(f)\) is exactly the canonical localization map
\(A_f\to A_g\).

% VI33-DETAILED-EXERCISE-03
\paragraph{Solution VI.33.E3.}
By definition, the stalk is the filtered colimit
\[
(\widetilde M)_{\mathfrak p}
=
\varinjlim_{\mathfrak p\in D(f)}\widetilde M(D(f))
=
\varinjlim_{f\notin\mathfrak p}M_f.
\]
Every fraction \(m/s\in M_{\mathfrak p}\) has \(s\notin\mathfrak p\), hence is represented in \(M_s\). Conversely,
the compatible maps \(M_f\to M_{\mathfrak p}\) for \(f\notin\mathfrak p\) induce a map from the colimit.  The two
constructions are inverse. Hence
\[
\boxed{(\widetilde M)_{\mathfrak p}\cong M_{\mathfrak p}}.
\]

% VI33-DETAILED-EXERCISE-04
\paragraph{Solution VI.33.E4.}
The affine module-sheaf formula gives
\[
\widetilde{A/I}(D(f))=(A/I)_f.
\]
Localization commutes with quotients:
\[
(A/I)_f
\cong
A_f/I A_f.
\]
If \(I_f\) denotes the extension \(IA_f\), then
\[
\boxed{\widetilde{A/I}(D(f))\cong A_f/I_f}.
\]
This is the coordinate ring of the intersection of the affine closed subscheme \(\Spec(A/I)\) with \(D(f)\).

% VI33-DETAILED-EXERCISE-05
\paragraph{Solution VI.33.E5.}
Restriction of an affine module sheaf to \(D(f)=\Spec A_f\) corresponds to localization of the module. Therefore
\[
\widetilde{A/I}|_{D(f)}=0
\iff
(A/I)_f=0.
\]
Equivalently, \(1=0\) in the localized quotient, so some power \(f^n\) maps to zero in \(A/I\):
\[
(A/I)_f=0
\iff
f^n\in I\text{ for some }n.
\]
Thus
\[
\boxed{\widetilde{A/I}|_{D(f)}=0\iff f\in\sqrt I}.
\]

% VI33-DETAILED-EXERCISE-06
\paragraph{Solution VI.33.E6.}
For a family of \(A\)-modules \((M_i)\), localization commutes with direct sums:
\[
\left(\bigoplus_iM_i\right)_f
\cong
\bigoplus_i(M_i)_f.
\]
Hence on every distinguished open
\[
\widetilde{\bigoplus_iM_i}(D(f))
\cong
\bigoplus_i\widetilde{M_i}(D(f)).
\]
These identifications are compatible with further localization, so they glue to
\[
\boxed{\widetilde{\bigoplus_iM_i}\cong\bigoplus_i\widetilde{M_i}}.
\]
In particular the tilde functor preserves finite and arbitrary direct sums.

% VI33-DETAILED-EXERCISE-07
\paragraph{Solution VI.33.E7.}
On \(D(f)\),
\[
\widetilde{M\otimes_A N}(D(f))
=
(M\otimes_A N)_f.
\]
Localization commutes with tensor product:
\[
(M\otimes_A N)_f
\cong
M_f\otimes_{A_f}N_f.
\]
The right side is exactly the section module of
\[
\widetilde M\otimes_{\mathcal O_X}\widetilde N
\]
on \(D(f)\).  Compatibility under restriction gives a sheaf isomorphism
\[
\boxed{
\widetilde{M\otimes_A N}
\cong
\widetilde M\otimes_{\mathcal O_X}\widetilde N.
}
\]

% VI33-DETAILED-EXERCISE-08
\paragraph{Solution VI.33.E8.}
For \(\mathfrak p\in\Spec A\),
\[
(\widetilde I)_{\mathfrak p}\cong I_{\mathfrak p}
\]
by the affine stalk formula.  Since localization preserves injections,
\[
I_{\mathfrak p}\hookrightarrow A_{\mathfrak p}.
\]
Thus the stalk is the localized ideal inside the local ring:
\[
\boxed{(\widetilde I)_{\mathfrak p}\cong I A_{\mathfrak p}\subseteq A_{\mathfrak p}}.
\]

% VI33-DETAILED-EXERCISE-09
\paragraph{Solution VI.33.E9.}
Start from the exact module sequence
\[
0\longrightarrow I\longrightarrow A\longrightarrow A/I\longrightarrow0.
\]
At each prime \(\mathfrak p\), localization is exact:
\[
0\longrightarrow I_{\mathfrak p}
\longrightarrow A_{\mathfrak p}
\longrightarrow (A/I)_{\mathfrak p}
\longrightarrow0.
\]
Using
\[
(\widetilde I)_{\mathfrak p}=I_{\mathfrak p},
\qquad
(\mathcal O_X)_{\mathfrak p}=A_{\mathfrak p},
\qquad
(\widetilde{A/I})_{\mathfrak p}=(A/I)_{\mathfrak p},
\]
the sheaf sequence is exact on every stalk. Therefore
\[
\boxed{
0\to\widetilde I\to\mathcal O_X\to\widetilde{A/I}\to0
}
\]
is exact.

% VI33-DETAILED-EXERCISE-10
\paragraph{Solution VI.33.E10.}
Let
\[
A=k[t],
\qquad
M=A/(t^3).
\]
The annihilator is
\[
\operatorname{Ann}_A(M)=(t^3).
\]
Since \(M\) is finitely generated,
\[
\operatorname{Supp}(\widetilde M)
=
V(t^3)
=
V(t).
\]
The prime ideals of \(k[t]\) containing \(t\) consist only of \((t)\). Hence
\[
\boxed{\operatorname{Supp}(\widetilde M)=\{(t)\}}.
\]
The exponent \(3\) records thickness in the module but is invisible in the underlying support.

% VI33-DETAILED-EXERCISE-11
\paragraph{Solution VI.33.E11.}
Set
\[
A=k[x,y],
\qquad
M=A/(x,y).
\]
After inverting \(x\), the element \(x\) becomes a unit but acts as zero on \(M\). Thus
\[
M_x=0.
\]
Similarly \(M_y=0\). Consequently
\[
\boxed{
\widetilde M|_{D(x)}=0,
\qquad
\widetilde M|_{D(y)}=0.
}
\]
This matches the support calculation: \(M\) is supported only at the origin \(V(x,y)\), which lies in neither
\(D(x)\) nor \(D(y)\).

% VI33-DETAILED-EXERCISE-12
\paragraph{Solution VI.33.E12.}
Let \(\mathcal E\) be finite free of rank \(r\) on a locally Noetherian scheme. On every affine open
\[
U=\Spec A
\]
over which it is free,
\[
\mathcal E|_U\cong\mathcal O_U^{\oplus r}\cong\widetilde{A^r}.
\]
The module \(A^r\) is finitely generated.  Under the locally Noetherian convention of VI/33, this is precisely the
local condition for coherence. Hence
\[
\boxed{\mathcal E\text{ is coherent}}.
\]

% VI33-DETAILED-EXERCISE-13
\paragraph{Solution VI.33.E13.}
Let \(M=(m_1,\ldots,m_r)\) be finitely generated.  A prime \(\mathfrak p\) is outside the support exactly when
\[
M_{\mathfrak p}=0.
\]
If \(\operatorname{Ann}M\not\subseteq\mathfrak p\), choose
\[
s\in\operatorname{Ann}M\setminus\mathfrak p.
\]
Then \(s\) is a unit in \(A_{\mathfrak p}\) and kills \(M_{\mathfrak p}\), so \(M_{\mathfrak p}=0\).

Conversely, suppose \(M_{\mathfrak p}=0\). For each generator \(m_i\), there exists
\(s_i\notin\mathfrak p\) with \(s_i m_i=0\). The product
\[
s=s_1\cdots s_r\notin\mathfrak p
\]
kills every generator, hence \(s\in\operatorname{Ann}M\). Therefore
\[
M_{\mathfrak p}\ne0
\iff
\operatorname{Ann}M\subseteq\mathfrak p.
\]
Thus
\[
\boxed{\operatorname{Supp}(\widetilde M)=V(\operatorname{Ann}M)}.
\]

% VI33-DETAILED-EXERCISE-14
\paragraph{Solution VI.33.E14.}
For
\[
M=A/(xy),
\qquad
A=k[x,y],
\]
the annihilator is \((xy)\). Hence
\[
\operatorname{Supp}(\widetilde M)
=
V(xy).
\]
Since
\[
V(xy)=V(x)\cup V(y),
\]
we get
\[
\boxed{\operatorname{Supp}(\widetilde{A/(xy)})=V(x)\cup V(y)}.
\]
The support is the union of the two coordinate axes.

% VI33-DETAILED-EXERCISE-15
\paragraph{Solution VI.33.E15.}
Let \(A\to B\) define
\[
f:\Spec B\to\Spec A.
\]
The pullback of \(\widetilde M\) is quasi-coherent. Its affine module is extension of scalars:
\[
B\otimes_A M.
\]
To see the formula locally, let \(g\in B\).  On \(D(g)\subseteq\Spec B\),
\[
(B\otimes_A M)_g
\cong
B_g\otimes_A M.
\]
This is exactly the module obtained from the sheaf-theoretic tensor definition of \(f^*\widetilde M\) on the
affine open \(D(g)\). Therefore
\[
\boxed{f^*\widetilde M\cong\widetilde{B\otimes_A M}}.
\]

% VI33-DETAILED-EXERCISE-16
\paragraph{Solution VI.33.E16.}
For the quotient map \(A\to A/I\),
\[
f^*\widetilde M
\cong
\widetilde{(A/I)\otimes_A M}.
\]
The standard tensor-quotient identity gives
\[
(A/I)\otimes_A M\cong M/IM.
\]
Hence
\[
\boxed{f^*\widetilde M\cong\widetilde{M/IM}}
\]
as an \(\mathcal O_{\Spec(A/I)}\)-module.

% VI33-DETAILED-EXERCISE-17
\paragraph{Solution VI.33.E17.}
Let \((U_i)\) be an affine open cover.  A sheaf is determined by its restrictions
\[
\mathcal F_i=\mathcal F|_{U_i}
\]
together with isomorphisms
\[
\varphi_{ij}:\mathcal F_i|_{U_i\cap U_j}
\longrightarrow
\mathcal F_j|_{U_i\cap U_j}
\]
satisfying the identity and cocycle conditions.  For a quasi-coherent sheaf each \(\mathcal F_i\) has the form
\(\widetilde{M_i}\). Ordinary sheaf gluing produces a unique \(\mathcal O_X\)-module from compatible data, and
quasi-coherence is local on \(X\). Thus compatible affine module-sheaf models determine \(\mathcal F\) uniquely up
to unique isomorphism.

% VI33-DETAILED-EXERCISE-18
\paragraph{Solution VI.33.E18.}
Let
\[
X=\Spec A.
\]
A sheaf morphism
\[
\varphi:\widetilde M\to\widetilde N
\]
induces on global sections an \(A\)-linear map
\[
u=\Gamma(X,\varphi):M\to N.
\]
On every distinguished open \(D(f)\), quasi-coherence forces the section map to be the localization
\[
u_f:M_f\to N_f.
\]
Therefore no further choices remain: \(\varphi\) is completely determined by \(u\).

% VI33-DETAILED-EXERCISE-19
\paragraph{Solution VI.33.E19.}
By E18, the map
\[
\operatorname{Hom}_{\mathcal O_X}(\widetilde M,\widetilde N)
\longrightarrow
\operatorname{Hom}_A(M,N),
\qquad
\varphi\mapsto\Gamma(X,\varphi),
\]
is injective. Conversely, any \(A\)-linear map \(u:M\to N\) localizes to compatible maps
\[
u_f:M_f\to N_f
\]
on all distinguished opens, and these define a sheaf morphism \(\widetilde u\). The constructions are inverse:
\[
\boxed{
\operatorname{Hom}_{\mathcal O_X}(\widetilde M,\widetilde N)
\cong
\operatorname{Hom}_A(M,N).
}
\]

% VI33-DETAILED-EXERCISE-20
\paragraph{Solution VI.33.E20.}
Let \(P\) be finitely generated projective. For every prime \(\mathfrak p\), the localization
\[
P_{\mathfrak p}
\]
is a finite free \(A_{\mathfrak p}\)-module. Choose a basis and lift its elements to \(P\).  Freeness and the basis
relations hold after inverting some element
\[
f\notin\mathfrak p.
\]
Thus
\[
P_f\cong A_f^{\oplus r}
\]
for some \(r\), and on \(D(f)\)
\[
\widetilde P|_{D(f)}
\cong
\widetilde{P_f}
\cong
\mathcal O_{D(f)}^{\oplus r}.
\]
Hence every prime has a distinguished trivializing neighborhood.

% VI33-DETAILED-EXERCISE-21
\paragraph{Solution VI.33.E21.}
Using compatibility of tilde with tensor products,
\[
\widetilde{A/(f)}\otimes\widetilde{A/(g)}
\cong
\widetilde{(A/(f))\otimes_A(A/(g))}.
\]
The tensor product of the quotient modules is
\[
(A/(f))\otimes_A(A/(g))
\cong
A/(f,g).
\]
Therefore
\[
\boxed{
\widetilde{A/(f)}\otimes\widetilde{A/(g)}
\cong
\widetilde{A/(f,g)}.
}
\]

% VI33-DETAILED-EXERCISE-22
\paragraph{Solution VI.33.E22.}
Let \(X=\Spec A\) and
\[
0\to\mathcal F'\to\mathcal F\to\mathcal F''\to0
\]
be exact with all three sheaves quasi-coherent. Under the affine equivalence write the sequence as
\[
0\to\widetilde{M'}\to\widetilde M\to\widetilde{M''}\to0.
\]
Exactness of tilde means the corresponding module sequence
\[
0\to M'\to M\to M''\to0
\]
is exact. Since
\[
\Gamma(X,\widetilde M)=M,
\]
taking global sections recovers this exact sequence, including surjectivity on the right.

For arbitrary sheaves on arbitrary spaces, global sections are only left exact: a section of the quotient sheaf
may exist locally and glue without admitting a global lift.  The affine quasi-coherent equivalence is precisely the
extra structure that removes this obstruction here.

% VI33-DETAILED-EXERCISE-23
\paragraph{Solution VI.33.E23.}
Work affine-locally on
\[
U=\Spec A
\]
with \(A\) Noetherian.  Let
\[
\mathcal G|_U\cong\widetilde M,
\qquad
\mathcal F|_U\cong\widetilde N
\]
where \(N\subseteq M\) are finitely generated \(A\)-modules. Then
\[
(\mathcal G/\mathcal F)|_U
\cong
\widetilde{M/N}.
\]
The quotient \(M/N\) is finitely generated. Hence the quotient sheaf is coherent on every such affine open, so
\[
\boxed{\mathcal G/\mathcal F\text{ is coherent}.}
\]

% VI33-DETAILED-EXERCISE-24
\paragraph{Solution VI.33.E24.}
On \(\Spec A\), the affine construction assigns
\[
D(f)\longmapsto M_f.
\]
For a graded ring \(S\), the standard opens of \(\Proj S\) are indexed by homogeneous elements \(f\), and their
coordinate rings are not the whole homogeneous localization \(S_f\), but its degree-zero part
\[
S_{(f)}=(S_f)_0.
\]
Accordingly, a graded module \(M\) should contribute
\[
M_{(f)}=(M_f)_0
\]
on \(D_+(f)\).  Thus the affine recipe survives, but homogeneous localization is followed by taking degree zero.

% VI33-DETAILED-EXERCISE-25
\paragraph{Solution VI.33.E25.}
Let
\[
0\to M'\to M\to M''\to0
\]
be exact. For every \(f\in A\), localization is an exact functor, so
\[
0\to M'_f\to M_f\to M''_f\to0
\]
is exact. By the distinguished-open formula,
\[
\widetilde{M'}(D(f))=M'_f,\quad
\widetilde M(D(f))=M_f,\quad
\widetilde{M''}(D(f))=M''_f.
\]
Thus the tilde sequence is exact on every distinguished open, and hence on every stalk.

This argument is special to these affine module sheaves and their localization description.  For a general short
exact sequence of sheaves, taking sections over an arbitrary open \(U\) need not preserve surjectivity on the
right; sheaf epimorphisms are locally surjective on germs, not necessarily globally surjective on sections.

% VI33-DETAILED-EXERCISE-26
\paragraph{Solution VI.33.E26.}
Let
\[
u:M\to N.
\]
The module sequences
\[
0\to\ker u\to M\to N
\]
and
\[
M\to N\to\operatorname{coker}u\to0
\]
are exact.  Exactness of the tilde functor gives
\[
0\to\widetilde{\ker u}\to\widetilde M\to\widetilde N
\]
and
\[
\widetilde M\to\widetilde N\to\widetilde{\operatorname{coker}u}\to0.
\]
By the universal characterization of kernels and cokernels in the sheaf category,
\[
\boxed{
\widetilde{\ker u}\cong\ker(\widetilde u),
\qquad
\widetilde{\operatorname{coker}u}\cong\operatorname{coker}(\widetilde u).
}
\]
Thus the affine module-sheaf equivalence respects the basic abelian-category operations.
